
\section{BCs in the form of time series}
The time series are stored in tables. For each BC code, there is one pair of (time , value) vectors. Thses vectors are filled with the time series provided in "hydro\_boundary\_time\_series.txt" at the init() stage. During the time loop, at a given simulation time $t$, 
%
%
\begin{enumerate}
%
\item for each BC code, compute the numerical values for all the BC types from the (time, values) vector pairs.
This is done with
\begin{verbatim}
for swBoundaryConditionTimeSeries* bc : bc_timeseries_list
{
    computeValBC(bc_time_series_list)
}
\end{verbatim}
%
\item for each boundary interface (\texttt{edge\_bnd[i]}, where i is the index), one must
\begin{enumerate}
\item get the BC type ("h", "z", "q", "fr", etc.). This is done by calling \texttt{edge\_bnd\_bc\_pointer[i]->getBCType}
\item get the numerical value. This is done by calling \texttt{edge\_bnd\_bc\_pointer[i]->getBCVal()}. The method carries out the interpolation from the (time, value) vector pairs
\item interpolate the numerical value for the BC from the (time, value) pair at time $t$.
\end{enumerate}
\end{enumerate}


%%%
\section{BCs in the form rating curves}

More to come in the next release. Stay tuned!